\documentclass[10pt,a4paper]{amsart}
\usepackage[margin=19mm]{geometry}
\usepackage{amsmath,amssymb,mathtools}
\usepackage[scr=rsfs,cal=euler]{mathalfa}
\usepackage{xfrac}
\usepackage[unicode,hidelinks,bookmarks=false]{hyperref}
\newcommand{\ie}{i.e.\ }
\newcommand{\cf}{cf.\ }
\newcommand{\resp}{respectively\ }
\newcommand{\Subsection}{\subsection}
\providecommand{\nobreakdash}{\nobreak\hbox{-}}
\setlength{\emergencystretch}{2em}
\multlinegap0pt
\usepackage{bm}
\usepackage{bbm}
\usepackage{dsfont}
\usepackage{xspace}
\usepackage{cancel}
\usepackage{mathtools}
\usepackage{amsthm}
\makeatletter
\@ifundefined{lemma}{\newtheorem{lemma}{Lemma}}{}
\@ifundefined{theorem}{\newtheorem{theorem}{Theorem}}{}
\makeatother
\title[Quantum Ergodicity and Thermalization in Interval Quantum Me]{Quantum Ergodicity and Thermalization in Interval Quantum Mechanics}
\author{Abbas Edalat}
\date{Source: arXiv:2606.00749v1 (CC BY 4.0)}
\begin{document}
\maketitle

\noindent\emph{Source and licence.} Quantum Ergodicity and Thermalization in Interval Quantum Mechanics. Version arXiv:2606.00749v1 (CC BY 4.0), distributed under the Creative Commons Attribution 4.0 International licence, as recorded in the arXiv licence metadata for this exact version and shown on the abstract page https://arxiv.org/abs/2606.00749v1. Full rights notice: licenses/arxiv-2606.00749-CC-BY-4.0.txt.\par\smallskip

\noindent\textit{Editorial scope.} Counted unit: the Theorem whose source label is \texttt{thm:basic} in the article (source0.tex lines 202--219, source label \texttt{thm:basic}), delivered together with the article's own complete original proof of it (source0.tex lines 221--257, one \texttt{proof} environment of 37 lines). No mathematics is written, repaired or re-derived: statement and proof are the article's own text line for line. Required context retained uncounted: lem:uniform (lines 151--166). Every definition, display and label of those lines is retained unchanged. Omitted from this package and not redistributed: 1--150, 167--201, 220--220, 258--372. Nothing inside a retained excerpt is omitted or altered: every retained body is delivered exactly as the article's lines are written, so no omission receipt of any kind accompanies this package. Citations: the retained excerpts carry no citation command at all, so no bibliography entry of the article is transcribed and no reference is redistributed. Reference adaptation: none was needed and none was made; every reference inside the delivered unit resolves inside it, so no unresolved reference or citation is produced. Printed slips kept literal: no printed word, symbol, hypothesis or number of the counted statement or of its proof was altered. Layout and class: the article's own class is \texttt{article} with packages including geometry, amsmath, amssymb, amsthm, braket, bm, booktabs, graphicx, hyperref; the wrapper is the lane's pinned \texttt{amsart} editorial wrapper at 10pt on A4, which restates the theorem-like environments the retained text needs and adds the article's own macro definitions that the retained text needs (none), with extra wrapper packages (bm, bbm, dsfont, xspace, cancel, mathtools). The delivered numbering is the unit's own. Assets: no retained span contains a figure environment, a graphics-inclusion command, a PDF-page inclusion, a table or a TikZ picture; no image file is copied into this package, the package declares no asset dependency and no image file is redistributed with it. Licence: version arXiv:2606.00749v1 of this article is distributed under the Creative Commons Attribution 4.0 International licence, as recorded in the arXiv licence metadata for this exact version and shown on the abstract page https://arxiv.org/abs/2606.00749v1. A full rights notice is delivered at licenses/arxiv-2606.00749-CC-BY-4.0.txt. Characters: every retained excerpt is pure ASCII, so the delivered engine, XeLaTeX with its default Latin Modern text font, reports no missing character.
\begin{lemma}[Uniform typicality over a parcel]
\label{lem:uniform}
Let \(O \subset \mathcal D(\mathcal H_\Delta)\) be a parcel with \(d_{\mathrm{eff}}(O)\ge D\). For any \(\varepsilon>0\) and any bounded observable \(A\), there exists a set \(\mathcal{G}\subset[0,\infty)\) (depending on \(O\) and \(\varepsilon\)) such that
\[
\limsup_{T\to\infty}\frac{1}{T}\,\lambda(\mathcal{G}^c\cap[0,T]) \le \min\left\{1,\; \frac{4N(\varepsilon) C(A)^2}{\varepsilon^2 D}\right\}
\]
and for all \(t\in\mathcal{G}\),
\[
\sup_{\rho\in O} \bigl|\operatorname{Tr}(\rho(t)A)-\operatorname{Tr}(\rho_{\mathrm{mc}}A)\bigr| \le \varepsilon.
\]
Here \(\lambda\) denotes Lebesgue measure on \([0,\infty)\),
\(\mathcal G^c=[0,\infty)\setminus\mathcal G\) denotes the complement
of \(\mathcal G\), and \(N(\varepsilon)\) is the number of centres in a
finite \(\varepsilon/(2\|A\|)\)-net covering the compact closure
\(\overline O\).
\end{lemma}

\begin{theorem}[Parcel thermalization under a uniform effective-dimension bound]
\label{thm:basic}
Let $O(0) \subset \mathcal D(\mathcal H_\Delta)$ be a parcel with \(d_{\mathrm{eff}}(O(0))\ge D\). For any fixed target precision \(\varepsilon>0\), let \(N(\varepsilon)\) be the finite covering number of \(\overline{O(0)}\) by trace-norm balls of radius \(\varepsilon/(2\|A\|)\). There exists a set \(\mathcal{G}\subset[0,\infty)\) such that
\[
\limsup_{T\to\infty}\frac{1}{T}\,\lambda(\mathcal{G}^c\cap[0,T]) \le \min\left\{1,\; \frac{4N(\varepsilon)C(A)^2}{\varepsilon^2 D}\right\}
\]
and for all \(t\in\mathcal{G}\) the following hold:

\begin{enumerate}
    \item \textbf{Thermalization of the possible set:} For any bounded observable \(A\), the expectation interval satisfies
    \[
    \operatorname{width}\bigl(E_{O(t)}(A)\bigr) \le 2\varepsilon, \qquad
    \bigl|\operatorname{centre}\bigl(E_{O(t)}(A)\bigr)-\operatorname{Tr}(\rho_{\mathrm{mc}}A)\bigr| \le \varepsilon.
    \]
    \item \textbf{Preservation of constants of motion:} If \([A,H]=0\), then \(E_{O(t)}(A)=E_{O(0)}(A)\) for all \(t\).
    \item \textbf{Uniformity of the bound:} For a fixed \(\varepsilon\), the asymptotic precision constraints in (1) depend only on \(\varepsilon\), not on the detailed internal shape of \(O(0)\).
\end{enumerate}
\end{theorem}

\begin{proof}
\textbf{Step 1 -- Apply uniform typicality to the parcel.}  
By Lemma~\ref{lem:uniform}, for the chosen fixed \(\varepsilon>0\), there exists a set \(\mathcal{G}\) such that for all \(t\in\mathcal{G}\) and all \(\rho\in O(0)\),
\[
\bigl|\operatorname{Tr}(\rho(t)A)-\operatorname{Tr}(\rho_{\mathrm{mc}}A)\bigr| \le \varepsilon. \tag{1}
\]

\textbf{Step 2 -- Derive the expectation interval bounds.}  
Recall that for any bounded observable \(A\),
\[
E_{O(t)}(A) = \bigl(\inf_{\rho\in O(0)}\operatorname{Tr}(\rho(t)A),\; \sup_{\rho\in O(0)}\operatorname{Tr}(\rho(t)A)\bigr).
\]
From (1), for all \(\rho\in O(0)\) and all \(t\in\mathcal{G}\),
\[
\operatorname{Tr}(\rho_{\mathrm{mc}}A)-\varepsilon \le \operatorname{Tr}(\rho(t)A) \le \operatorname{Tr}(\rho_{\mathrm{mc}}A)+\varepsilon.
\]
Taking the infimum and supremum over \(\rho\in O(0)\) yields
\[
\operatorname{Tr}(\rho_{\mathrm{mc}}A)-\varepsilon \le \inf_{\rho\in O(0)}\operatorname{Tr}(\rho(t)A) \le \sup_{\rho\in O(0)}\operatorname{Tr}(\rho(t)A) \le \operatorname{Tr}(\rho_{\mathrm{mc}}A)+\varepsilon.
\]
Hence
\[
\bigl|\operatorname{centre}(E_{O(t)}(A))-\operatorname{Tr}(\rho_{\mathrm{mc}}A)\bigr| \le \varepsilon,\qquad
\operatorname{width}(E_{O(t)}(A)) \le 2\varepsilon.
\]

\textbf{Step 3 -- Constants of motion.}  
If \([A,H]=0\), then \(U(t)^\dagger A U(t)=A\) for all \(t\). Hence for any \(\rho\in O(0)\),
\[
\operatorname{Tr}(\rho(t)A) = \operatorname{Tr}\bigl(U(t)\rho U(t)^\dagger A\bigr)
= \operatorname{Tr}\bigl(\rho U(t)^\dagger A U(t)\bigr) = \operatorname{Tr}(\rho A).
\]
Thus \(E_{O(t)}(A)=E_{O(0)}(A)\) for all \(t\).

\textbf{Step 4 -- Uniformity of the bound.}  
The structural bounds depend only on \(\varepsilon\). The good time set \(\mathcal{G}\) depends on \(O(0)\) through the geometry of its finite cover, but the existence of a valid macro-set is sufficient; the uniformity refers explicitly to the precision of the late-time expectation intervals, which is independent of the initial configuration. 
\end{proof}

\end{document}
